% !PRINTED 130
\weberpage{130}
%<<CONT
\gesperrt{divisor of $\mu$, then there are $\nu$ and no more elements in $P$ whose degree divides $\nu$.}

If we apply this theorem to $\nu = t$, it follows that apart from $T$ there can be no elements in $P$ whose degree $= t$. But if $c$ is any element of $P$, then the group $c^{-1}Tc$ conjugate to $T$ is likewise of degree $t$, and must therefore be identical with $T$. It is thereby proved that $T$ is a normal divisor of $P$.

For the proof of theorem $\alpha$) one applies complete induction. It is evidently correct when $\mu = 1$, $\nu = n$, and we therefore assume as proved:

$\alpha'$) \gesperrt{If $\mu'\nu'$ is divisible by no square, every prime divisor of $\nu'$ greater than every prime divisor of $\mu'$, and the number of prime divisors of $\mu'$ smaller than the number of prime divisors of $\mu$, then in every group of degree $\mu'\nu'$ there are exactly $\nu'$ elements whose degree is a divisor of $\nu'$.}

To derive from this the proof of theorem $\alpha$), we denote by $p$ the greatest prime divisor of $\mu$ and by $Q$ a group of degree $p$ contained in $P$, whose existence is established by the \gesperrt{Cauchy-Sylow} theorem. Then all prime factors of $\nu$ are greater than $p$. As in \S.\,30 we denote by $R$ the group which consists of all elements $c$ of $P$ satisfying the condition
\[
c^{-1}Qc = Q
\]
whose degree we denote, as above, by $pr$; $j$ signifies, as before, the index of the divisor $R$ of $P$. We set $r = r_1 r_2$, and understand by $r_1$ the greatest common divisor of $\mu$ and $r$, so that $r_2$ is the greatest common divisor of $\nu$ and $r$. Then
\[
n = \mu\nu = r_1 p r_2 j.
\]
By the hypothesis $\alpha'$), $P$ now contains exactly $p\nu$ elements whose degree is a divisor of $p\nu$; let the totality of these elements be denoted by $U$. But the group $R$ contains, likewise by $\alpha'$), exactly $pr_2$ elements whose degree is a divisor of $pr_2$, and hence a divisor of $p\nu$. This system we shall denote by $V$. Evidently all elements of $V$ are at the same time contained in $U$.
